\medskip

On considère la suite numérique $\left(u_{n}\right)$ définie par son premier terme $u_{0}=2$ et pour tout entier naturel $n$, par :


\[u_{n+1}=\dfrac{2 u_{n}+1}{u_{n}+2}\]

On admet que la suite $\left(u_{n}\right)$ est bien définie.

\begin{enumerate}
\item Calculer le terme $u_{1}$.
\item On définit la suite $\left(a_{n}\right)$ pour tout entier naturel $n$, par :

\[a_{n}=\dfrac{u_{n}}{u_{n}-1}\]

On admet que la suite $\left(a_{n}\right)$ est bien définie.

	\begin{enumerate}
		\item Calculer $a_{0}$ et $a_{1}$.

		\item Démontrer que, pour tout entier naturel $n$, $a_{n+1}=3 a_{n}-1$.

		\item Démontrer par récurrence que, pour tout entier naturel $n$ supérieur ou égal à~1,

\[a_{n} \geqslant 3 n-1\]

		\item En déduire la limite de la suite $\left(a_{n}\right)$.
	\end{enumerate}
\item On souhaite étudier la limite de la suite $\left(u_{n}\right)$.
	\begin{enumerate}
		\item Démontrer que pour tout entier naturel $n$, $u_{n}=\dfrac{a_{n}}{a_{n}-1}$.

		\item En déduire la limite de la suite $\left(u_{n}\right)$.
		\end{enumerate}
\item On admet que la suite $\left(u_{n}\right)$ est décroissante.

On considère le programme suivant écrit en langage Python :

\begin{center}
\begin{ttfamily}
\begin{tabularx}{8cm}{|l @{\quad} X|}\hline
	1& \textbf{def} algo(p):\\
	2& \quad u=2\\
	3& \quad n=0\\
	4& \quad \textbf{while} u-1>p:\\
	5& \quad\quad u=(2*u+1)/(u+2)\\
	6& \quad\quad n=n+1\\
	7& \quad \textbf{return} (n,u)\\ \hline
\end{tabularx}
\end{ttfamily}
\end{center}

	\begin{enumerate}
		\item Interpréter les valeurs \texttt{n} et \texttt{u} renvoyées par l'appel de la fonction \texttt{algo(p)} dans le contexte de l'exercice.
		\item Donner, sans justifier, la valeur de $n$ pour $p=0,001$.
	\end{enumerate}
\end{enumerate}
